\documentclass[12pt]{article}

\usepackage[utf8]{inputenc}
\usepackage[uzbek]{babel}
\usepackage{amsmath, amssymb, amsthm}
\usepackage{geometry}
\usepackage{graphicx}
\usepackage{tikz}
\usepackage{booktabs}
\usepackage{caption}
\usepackage{fontawesome5}
\usepackage{tcolorbox}

\addto\captionsuzbek{\renewcommand{\tablename}{Jadval}}

\geometry{a4paper, margin=1in}

\title{Fazzi to'plamlar}
\author{Shahrizoda Kenjaboyeva}
\date{\today}

\begin{document}

\maketitle

\section{Fazzi to'plam tushunchasi}

$(\mu_A(x), x)$ ko'rinishidagi juftliklar to'plamiga \textbf{A fazzi to'plami} deyiladi.
Bu yerda $x \in X$ va $\mu_A(x)$-$x$ ning $A$ to'plamga a'zolik funksiyasi.

\[\mu_A(X): X \to [0, 1]\]

Agar $B \subseteq X$ bo'lsa, uning a'zolik funksiyasi quyidagicha aniqlanadi:
\[\mu_B(x)=
\begin{cases}
1, & x\in B \\
0, & x\notin B
\end{cases}\]
Bundan ko'rinadiki oddiy to'plamlar fazzi to'plamlarning xususiy holidir.

\subsection{Asosiy tushunchalar}

\begin{itemize}
	\item \textbf{Tashuvchi:} $\mu_A(x)>0$ bo'lgan barcha $x$ lar $A$ ning tashuvchilari deyiladi.
	\item \textbf{Bir nuqtali fazzi to'plam:} Tashuvchisi faqat bitta elementdan iborat bo'lgan fazzi to'plam: $A=\frac{\mu}{x}$.
	\item \textbf{O'tish nuqtasi:} $\mu_A(x)=0.5$ bo'lgan $x$ nuqta.
\end{itemize}

\section{Fazzi to'plamlarni ifodalash}

\begin{itemize}
	\item $A=\int \frac{\mu_A(x)}{x} \, dx$
	
Tashuvchilari chekli elementlardan iborat $A$ fazzi to'plamni quyidagicha belgilash ham mumkin:

	\item $A=\frac{\mu_1}{x_1}+\frac{\mu_2}{x_2}+...+\frac{\mu_n}{x_n}$
	$A=\sum_{i=1}^{n}\frac{\mu_i}{x_i}$
	
\begin{table}[h]
	\centering
	\caption{Fazzi to'plamlarni ifodalashning yana bir usuli:}
	\vspace{2mm}
	\begin{tabular}{c|c|c|c|c}
		\toprule
		$\mu_i$ & $\mu_1$ & $\mu_2$ & ...& $\mu_n$ \\ \midrule
		$x_i$ & $x_1$ & $x_2$ & ... &$x_n$ \\ \bottomrule
	\end{tabular}
\end{table}
	
\end{itemize}

\begin{quote}
	\faExclamationTriangle\ \textbf{Eslatma:} Bularning hech biri integral, summa yoki taqsimot jadvali emas!
\end{quote}


\end{document}



